\documentclass[12pt]{article}
% Préambule commun à tous les documents extraits de « La Revue CAPES »
\usepackage[T1]{fontenc}
\usepackage[utf8]{inputenc}
\usepackage{lmodern}
\usepackage{amsmath,amssymb,amsthm,amsfonts,mathrsfs}
\usepackage{setspace}
\usepackage{listings}
\usepackage{pgfplots}
\pgfplotsset{compat=1.18}
\usepackage{longtable}
\usepackage{adjustbox}
\usepackage{lettrine}
\usepackage{tkz-tab}
\usepackage{changepage}
\usepackage{pdflscape}
\usepackage{graphicx}
\usepackage{tikz}
\usepackage{xcolor}
\usepackage[a4paper,top=2cm,bottom=2.5cm,left=2.5cm,right=2cm]{geometry}
\usepackage{fancyhdr}
\usepackage{draftwatermark}
\SetWatermarkText{La RevueCapes}
\SetWatermarkLightness{0.9}
\SetWatermarkScale{2}
\usepackage{hyperref}
\hypersetup{colorlinks=true,linkcolor=blue!50!black,urlcolor=blue!50!black}

\definecolor{revue}{RGB}{20,60,120}

\pagestyle{fancy}
\fancyhf{}
\renewcommand{\headrulewidth}{0.4pt}
\setlength{\headheight}{15pt}

% \entete{Matière}{Titre}{Type de document}
\newcommand{\entete}[3]{%
  \fancyhead[L]{\small\textsc{La Revue CAPES} -- #1}%
  \fancyhead[R]{\small #2}%
  \fancyfoot[C]{\thepage}%
  \thispagestyle{plain}%
  \begin{center}
    {\color{revue}\rule{\textwidth}{1.2pt}}\\[4pt]
    {\small\textsc{Concours directs CAP-CEG / CAPES -- Burkina Faso}}\\[6pt]
    {\LARGE\bfseries\color{revue} #2}\\[6pt]
    {\large\itshape #1 -- #3}\\[2pt]
    {\color{revue}\rule{\textwidth}{1.2pt}}
  \end{center}
  \vspace{0.5cm}%
}

% Encadré pour les remarques ajoutées dans les corrigés
\newcommand{\remarque}[1]{\par\medskip\noindent\fcolorbox{revue}{revue!6}{\parbox{\dimexpr\linewidth-2\fboxsep-2\fboxrule}{\textbf{\color{revue}Remarque.} #1}}\par\medskip}


\begin{document}
\entete{Culture générale}{Corrigé du sujet type de dissertation n°3}{Correction}
\begin{spacing}{1.5}

\textbf{ Introduction}

\textbf{Contexte(préambule) :} L'école est souvent perçue comme le chemin principal vers le succès, en offrant des compétences, des diplômes et des opportunités d'emploi. Cependant, de nombreuses réussites dans la vie ont montré que l'école n'est pas le seul moyen d'atteindre ses objectifs.

\textbf{ Problématique :} Dans quelle mesure l'école est-elle indispensable pour réussir ? Existe-t-il d'autres voies permettant d'accomplir ses ambitions et d'obtenir une reconnaissance sociale ou personnelle ?

\textbf{ Annonce du plan :} Nous discuterons tout d'abord de l'importance de l'école dans la réussite, puis des limites de cette institution, avant d'explorer d'autres moyens de réussir dans la vie.

\quad
I. L'école : un moyen structuré et efficace pour réussir
L'école reste un pilier fondamental pour construire une base solide en vue du succès.

1. Accès aux connaissances et aux compétences
L'école fournit des connaissances théoriques et pratiques indispensables dans de nombreux domaines professionnels.

Exemple : Des métiers comme médecin, ingénieur ou avocat nécessitent des diplômes spécifiques.

2. Reconnaissance sociale
Le système scolaire attribue des diplômes qui sont des critères de validation universels dans le monde du travail.

Exemple : Un employeur privilégiera souvent un candidat ayant un diplôme pertinent.

3. Développement personnel et social
L'école permet de développer des compétences comme la discipline, la collaboration, et la pensée critique.

Exemple : Les activités parascolaires (clubs, compétitions) renforcent les aptitudes interpersonnelles.

\quad
II. Les limites de l'école comme unique moyen de réussite
Malgré ses avantages, l'école ne garantit pas toujours le succès et peut ne pas convenir à tout le monde.

1. L'école ne développe pas toujours les talents individuels.

Le système scolaire, souvent standardisé, peut négliger les talents et intérêts spécifiques des individus.

Exemple : De nombreux artistes, sportifs ou entrepreneurs brillants ont échoué à l'école parce que leur potentiel n'y était pas valorisé.

2. Inadéquation avec le marché du travail

Les diplômes scolaires ne garantissent pas toujours un emploi, en particulier dans des contextes où le chômage est élevé.

Exemple : Un diplômé sans expérience pratique peut avoir des difficultés à trouver un emploi adapté.

3. L'école ne prend pas en compte l'intelligence pratique
Certaines formes d'intelligence, comme l'intelligence émotionnelle, sociale ou manuelle, ne sont pas suffisamment valorisées dans le système scolaire.

Exemple : Les artisans ou commerçants réussissent souvent sans avoir fait d'études supérieures.

\quad
III. D'autres moyens de réussir dans la vie
Le succès peut être atteint par d'autres voies, notamment grâce à la persévérance, la créativité et l'adaptabilité.

1. L'apprentissage informel et autodidacte

L'ère numérique offre une multitude de ressources éducatives gratuites ou accessibles, permettant d'apprendre sans passer par l'école traditionnelle.

Exemple : Bill Gates, autodidacte dans l'informatique, a quitté l'université mais a fondé Microsoft.

2. Les métiers manuels et techniques

Des métiers comme maçon, couturier ou mécanicien ne nécessitent pas de longues études mais peuvent mener à un succès économique et personnel.

Exemple : De nombreux entrepreneurs locaux ont bâti des entreprises prospères en s'appuyant sur leurs compétences pratiques.

3. La créativité et l'innovation

Les secteurs artistiques (musique, cinéma, mode) montrent que des individus peuvent réussir grâce à leur talent et leur originalité, sans passer par un cursus scolaire classique.

Exemple : Des artistes comme Bob Marley ou des créateurs comme Coco Chanel n'ont pas suivi de formation académique rigoureuse.

4. Les réseaux sociaux et le marketing personnel
À l'ère d'Internet, les plateformes numériques permettent à des individus de se faire connaître et de réussir, indépendamment de leur parcours scolaire.

Exemple : Les influenceurs et entrepreneurs numériques qui bâtissent leur carrière sur des compétences pratiques et du marketing personnel.

\textbf{ Conclusion}

\textbf{ Synthèse }: L'école constitue un moyen structuré et efficace de réussir, mais elle n'est pas le seul chemin. Le succès peut également être atteint par des voies alternatives, comme l'apprentissage autodidacte, les métiers techniques ou la créativité.

\textbf{Réflexion personnelle :} L'essentiel pour réussir réside dans la capacité à identifier ses talents, à persévérer face aux obstacles et à s'adapter à un monde en constante évolution.

\textbf{ Ouverture :} Plutôt que d'opposer école et autres moyens de réussite, il serait pertinent de concevoir un système éducatif plus inclusif et flexible, capable de valoriser les divers talents et aspirations des individus.






\quad

\end{spacing}
\end{document}
